\chapter{Introduction}
Recent progress in hardware technology allowed computers to store an increasing amount of large data.
A massive amount of data is produced and consumed each moment on the internet.
In 2021, 2.5 quintillion bytes were created every day.
A lot of this data is even publicly available, as reported by Archive.org in 2012, when they archived ten petabytes in total.
An increase in data and surrounding importance makes it also important to teach about \emph{correct} data visualization.

\section{Why Visualize Data?}
Data visualizations are vitally important to extract important information from often clouded data.
Statistical measures such as mean, variance, and correlation are often not conclusive enough to do this job\footnote{\url{https://jumpingrivers.github.io/datasauRus/}}.

The human visual system is a comparatively high-bandwidth link at 10 million bits per second to the brain, therefore it can be used to convey a lot of information very fast.

\subsection{History of Visualization}
Data visualizations introduced different goals.
At first, they were only used to convey already known information---the goal was to underline and present the facts.
Then, they were also used to confirm or reject an already set hypothesis.
This is called \emph{confirmative analysis}.
Finally, they are also used to find a hypothesis from interesting data.
The latter is called \emph{explorative analysis}.

The standard pipeline from data to visualization also changed.
During the 1990s, this pipeline was straightforward: First data is filtered and then rendered.
Recently, user interability became more important, leading to interactive diagrams, where each step from data to diagram can be modified in-time.
Today, due to the popularity of machine learning algorithms such as transformers, a recent trend in visualization is also to use machine learning to visualize data.
This method was conceptualized as machine learning models are usually hard to understand for a human.
Therefore, a new goal is to use machine learning to create a model which solves the same task, but is humanly understandable.
